
\documentclass{article}
\usepackage{colortbl}
\usepackage{makecell}
\usepackage{multirow}
\usepackage{supertabular}

\begin{document}

\newcounter{utterance}

\centering \large Interaction Transcript for game `cryptolect', experiment `basic', episode 3 with LLP: Large Language Problems\_\_qwen3.5{-}9b\_\_llp{-}final.
\vspace{24pt}

{ \footnotesize  \setcounter{utterance}{1}
\setlength{\tabcolsep}{0pt}
\begin{supertabular}{c@{$\;$}|p{.15\linewidth}@{}p{.15\linewidth}p{.15\linewidth}p{.15\linewidth}p{.15\linewidth}p{.15\linewidth}}
    \# & \multicolumn{2}{c}{Player} && \multicolumn{2}{c}{Game Master} \\
    \hline

    \theutterance \stepcounter{utterance}  
    & & & \multicolumn{4}{p{0.6\linewidth}}{
        \cellcolor[rgb]{0.9,0.9,0.9}{
            \makecell[{{p{\linewidth}}}]{
                \texttt{\tiny{[P1$\langle$GM]}}
                \texttt{You are a field linguist. You have made contact with the only speaker of an unknown} \\
\texttt{language (the "informant"). Nobody has ever written down this language before. Your task} \\
\texttt{is to work out how it works and then translate sentences into it correctly.} \\
\\ 
\texttt{You will study the language by asking the informant to translate English sentences for} \\
\texttt{you. You may only talk about a small fixed world, using exactly these English words:} \\
\\ 
\texttt{Nouns: boat, ship, hill, tree, stone, dog, farmer, road} \\
\texttt{Adjectives: (this language has no adjectives)} \\
\texttt{Verbs (write them in this base form): watch, hold, carry, hear, chase} \\
\texttt{Articles: the (definite), a (indefinite)} \\
\texttt{Number words: two, three (no number word = one / singular)} \\
\texttt{These are ALL the words in play {-} the language has no other words.} \\
\\ 
\texttt{HOW TO WRITE AN ENGLISH SENTENCE} \\
\texttt{Every sentence is one clause with a subject, a verb, and an object, written in this fixed English notation:} \\
\texttt{    SUBJECT  verb  OBJECT} \\
\texttt{Each of SUBJECT and OBJECT is a noun phrase of the form:} \\
\texttt{    {[}the | a{]}  {[}number word{]}  {[}adjective ...{]}  noun} \\
\texttt{{-} 'the' = definite, 'a' = indefinite (singular only); leave the article off for an indefinite plural.} \\
\texttt{{-} a number word (two, three) sets the count; with no number word the noun is singular.} \\
\texttt{Examples of well{-}formed English sentences you may send:} \\
\texttt{    a boat watch a ship} \\
\texttt{    three boats watch a ship} \\
\\ 
\texttt{WHAT YOU ALREADY KNOW} \\
\texttt{The informant has already translated these sentences for you. Each line is} \\
\texttt{"<English>  =>  <the language>":} \\
\texttt{    three trees watch two boats  =>  dugunu nu deniru gutinu ga} \\
\texttt{    three dogs hear a dog  =>  tipi puga tipinu ga} \\
\texttt{    three dogs carry three ships  =>  pirenu ga geti tipinu ga} \\
\texttt{    two dogs chase two boats  =>  dugunu nu buro tipinu nu} \\
\texttt{    two roads hold three trees  =>  gutinu ga nira dobudinu nu} \\
\texttt{    three farmers watch three farmers  =>  nidinu ga deniru nidinu ga} \\
\texttt{    three stones carry two roads  =>  dobudinu nu geti bodenu ga} \\
\texttt{    two trees hold three hills  =>  ruditinu ga nira gutinu nu} \\
\\ 
\texttt{HOW THE SESSION WORKS} \\
\texttt{1. INVESTIGATION: You have 4 rounds of questioning. In each round you may ask} \\
\texttt{   the informant to translate up to 3 English sentences (so at most} \\
\texttt{   12 probes in total). The informant answers each one truthfully and} \\
\texttt{   exactly. Use these answers to figure out the words and the grammar (word order, how} \\
\texttt{   number/case/definiteness are marked, how words must agree, and any sound changes).} \\
\texttt{   Choose your probes deliberately {-} test specific hypotheses.} \\
\texttt{2. EXAM: After the investigation rounds, the informant will give you 10 NEW} \\
\texttt{   English sentences. You must translate each one INTO the language. Every answer is} \\
\texttt{   graded by EXACT match against how the informant would say it {-} every word, ending and} \\
\texttt{   space must be right (capitalisation and surrounding spaces are ignored). The exam} \\
\texttt{   sentences will use the same vocabulary but in combinations you have not been shown, so} \\
\texttt{   you must induce the general rules, not just memorise examples.} \\
\\ 
\texttt{RESPONSE FORMAT (for every investigation round)} \\
\texttt{List your probe sentences under a 'PROBES:' line, one per line. A THOUGHTS line is} \\
\texttt{optional {-} keep any reasoning brief so you always finish writing your PROBES:} \\
\texttt{THOUGHTS: <optional; brief>} \\
\texttt{PROBES:} \\
\texttt{1. <an English sentence, written as described above>} \\
\texttt{2. <another English sentence>} \\
\texttt{(up to 3 lines; at least one)} \\
\\ 
\texttt{Begin your first round of probes now.} \\
            }
        }
    }
     \\ \\

    \theutterance \stepcounter{utterance}  
    & \multicolumn{4}{p{0.6\linewidth}}{
        \cellcolor[rgb]{0.561,0.737,0.561}{
            \makecell[{{p{\linewidth}}}]{
                \texttt{\tiny{[P1$\rangle$GM]}}
                \texttt{THOUGHTS: I need to analyze the given translations to understand the grammar. Let me break down what I know:} \\
\\ 
\texttt{From the examples:} \\
\texttt{{-} "three trees watch two boats" => "dugunu nu deniru gutinu ga"} \\
\texttt{{-} "three dogs hear a dog" => "tipi puga tipinu ga"} \\
\texttt{{-} "three dogs carry three ships" => "pirenu ga geti tipinu ga"} \\
\texttt{{-} "two dogs chase two boats" => "dugunu nu buro tipinu nu"} \\
\texttt{{-} "two roads hold three trees" => "gutinuga nira dobudinu nu"} \\
\texttt{{-} "three farmers watch three farmers" => "nidinuga deniru nidinuga"} \\
\texttt{{-} "three stones carry two roads" => "dobudinu nu geti bodenu ga"} \\
\texttt{{-} "two trees hold three hills" => "ruditinuga nira gutinu nu"} \\
\\ 
\texttt{Observations:} \\
\texttt{1. Word order seems to be: OBJECT (with verb) then SUBJECT (with verb)} \\
\texttt{2. The verb appears to be in a specific form {-} often ending in {-}iru or {-}inu} \\
\texttt{3. "ga" appears at the end of some sentences, "nu" at the end of others} \\
\texttt{4. "nira" appears in the middle of sentences} \\
\texttt{5. The first noun in each pair seems to be the object, second is the subject} \\
\\ 
\texttt{Let me test hypotheses about:} \\
\texttt{{-} What "ga" and "nu" mark (subject/object?)} \\
\texttt{{-} What "nira" does} \\
\texttt{{-} How verbs are formed} \\
\texttt{{-} How number is marked} \\
\\ 
\texttt{PROBES:} \\
\texttt{1. a boat watch a tree} \\
\texttt{2. two ships chase a dog} \\
\texttt{3. a farmer hold a stone} \\
            }
        }
    }
    & & \\ \\

    \theutterance \stepcounter{utterance}  
    & & & \multicolumn{2}{p{0.3\linewidth}}{
        \cellcolor[rgb]{0.9,0.9,0.9}{
            \makecell[{{p{\linewidth}}}]{
                \texttt{\tiny{[GM$|$GM]}}
                {'kind': 'investigation', 'round': 0, 'probes_issued': 3, 'malformed': 0, 'probes': ['a boat watch a tree', 'two ships chase a dog', 'a farmer hold a stone']}
            }
        }
    }
    & & \\ \\

    \theutterance \stepcounter{utterance}  
    & & & \multicolumn{4}{p{0.6\linewidth}}{
        \cellcolor[rgb]{0.9,0.9,0.9}{
            \makecell[{{p{\linewidth}}}]{
                \texttt{\tiny{[P1$\langle$GM]}}
                \texttt{TRANSLATIONS:} \\
\texttt{1. "a boat watch a tree"  =>  guti deniru dugu} \\
\texttt{2. "two ships chase a dog"  =>  tipi buro pirenu nu} \\
\texttt{3. "a farmer hold a stone"  =>  bode nira nidi} \\
\texttt{STATUS: investigation round 1 of 4 done. Probes used: 3 of 12. 3 investigation round(s) left before the exam (10 sentences).} \\
\texttt{Send your next probes now (a PROBES list; brief THOUGHTS optional).} \\
            }
        }
    }
     \\ \\

    \theutterance \stepcounter{utterance}  
    & \multicolumn{4}{p{0.6\linewidth}}{
        \cellcolor[rgb]{0.561,0.737,0.561}{
            \makecell[{{p{\linewidth}}}]{
                \texttt{\tiny{[P1$\rangle$GM]}}
                \texttt{THOUGHTS: From the translations, I can see:} \\
\texttt{{-} "guti deniru dugu" = "a boat watch a tree" {-} so "guti" is boat, "deniru" is watch, "dugu" is tree} \\
\texttt{{-} "tipi buro pirenu nu" = "two ships chase a dog" {-} so "tipi" is dog, "buro" is chase, "pirenu" is ships (plural), "nu" at end} \\
\texttt{{-} "bode nira nidi" = "a farmer hold a stone" {-} so "bode" is stone, "nira" is hold, "nidi" is farmer} \\
\\ 
\texttt{Comparing with earlier data:} \\
\texttt{{-} "dugunu nu deniru gutinu ga" = "three trees watch two boats" {-} "dugunu" is trees (plural), "deniru" is watch, "gutinuga" is boats (plural)} \\
\texttt{{-} "tipi puga tipinu ga" = "three dogs hear a dog" {-} "tipi" is dog, "puga" is hear, "tipinu" is dogs (plural)} \\
\texttt{{-} "gutinuga nira dobudinu nu" = "two roads hold three trees" {-} "gutinuga" is roads (plural), "nira" is hold, "dobudinu" is trees (plural)} \\
\\ 
\texttt{So "nira" = hold, "deniru" = watch, "buro" = chase, "puga" = hear, "geti" = carry} \\
\\ 
\texttt{Pattern: The verb seems to be in the form where the first noun is the object and second is the subject. The verb form changes based on what it's acting on.} \\
\\ 
\texttt{Let me test more to understand the verb forms and number marking.} \\
\\ 
\texttt{PROBES:} \\
\texttt{1. a dog chase a boat} \\
\texttt{2. two trees hold a road} \\
\texttt{3. three stones watch a farmer} \\
            }
        }
    }
    & & \\ \\

    \theutterance \stepcounter{utterance}  
    & & & \multicolumn{2}{p{0.3\linewidth}}{
        \cellcolor[rgb]{0.9,0.9,0.9}{
            \makecell[{{p{\linewidth}}}]{
                \texttt{\tiny{[GM$|$GM]}}
                {'kind': 'investigation', 'round': 1, 'probes_issued': 3, 'malformed': 0, 'probes': ['a dog chase a boat', 'two trees hold a road', 'three stones watch a farmer']}
            }
        }
    }
    & & \\ \\

    \theutterance \stepcounter{utterance}  
    & & & \multicolumn{4}{p{0.6\linewidth}}{
        \cellcolor[rgb]{0.9,0.9,0.9}{
            \makecell[{{p{\linewidth}}}]{
                \texttt{\tiny{[P1$\langle$GM]}}
                \texttt{TRANSLATIONS:} \\
\texttt{1. "a dog chase a boat"  =>  dugu buro tipi} \\
\texttt{2. "two trees hold a road"  =>  dobudi nira gutinu nu} \\
\texttt{3. "three stones watch a farmer"  =>  nidi deniru bodenu ga} \\
\texttt{STATUS: investigation round 2 of 4 done. Probes used: 6 of 12. 2 investigation round(s) left before the exam (10 sentences).} \\
\texttt{Send your next probes now (a PROBES list; brief THOUGHTS optional).} \\
            }
        }
    }
     \\ \\

    \theutterance \stepcounter{utterance}  
    & \multicolumn{4}{p{0.6\linewidth}}{
        \cellcolor[rgb]{0.561,0.737,0.561}{
            \makecell[{{p{\linewidth}}}]{
                \texttt{\tiny{[P1$\rangle$GM]}}
                \texttt{THOUGHTS: Now I'm seeing clearer patterns:} \\
\\ 
\texttt{1. "dugu buro tipi" = "a dog chase a boat" {-} "dugu" is dog, "buro" is chase, "tipi" is boat} \\
\texttt{2. "dobudi nira gutinu nu" = "two trees hold a road" {-} "dobudi" is trees (plural), "nira" is hold, "gutinuga" is road (singular)} \\
\texttt{3. "nidi deniru bodenu ga" = "three stones watch a farmer" {-} "nidi" is stones (plural), "deniru" is watch, "bodenu" is farmer (singular)} \\
\\ 
\texttt{So the pattern is: OBJECT{-}VERB{-}SUBJECT} \\
\texttt{{-} Object comes first, then verb, then subject} \\
\texttt{{-} The verb form changes based on what it's acting on} \\
\\ 
\texttt{Number marking:} \\
\texttt{{-} Singular nouns: dugu, tipi, gutinu, nidi, bodenu} \\
\texttt{{-} Plural nouns: dugunu, dobudi, bodenu (wait, bodenu is singular farmer, but bodenu ga is plural farmers)} \\
\\ 
\texttt{Actually looking at "bodenu ga" = farmers (plural) vs "nidi" = stones (plural), so "ga" marks plural at the end of the subject.} \\
\\ 
\texttt{Let me test more to confirm the verb forms and understand how they change.} \\
\\ 
\texttt{PROBES:} \\
\texttt{1. a road chase a dog} \\
\texttt{2. a tree hold a boat} \\
\texttt{3. a farmer chase a tree} \\
            }
        }
    }
    & & \\ \\

    \theutterance \stepcounter{utterance}  
    & & & \multicolumn{2}{p{0.3\linewidth}}{
        \cellcolor[rgb]{0.9,0.9,0.9}{
            \makecell[{{p{\linewidth}}}]{
                \texttt{\tiny{[GM$|$GM]}}
                {'kind': 'investigation', 'round': 2, 'probes_issued': 3, 'malformed': 0, 'probes': ['a road chase a dog', 'a tree hold a boat', 'a farmer chase a tree']}
            }
        }
    }
    & & \\ \\

    \theutterance \stepcounter{utterance}  
    & & & \multicolumn{4}{p{0.6\linewidth}}{
        \cellcolor[rgb]{0.9,0.9,0.9}{
            \makecell[{{p{\linewidth}}}]{
                \texttt{\tiny{[P1$\langle$GM]}}
                \texttt{TRANSLATIONS:} \\
\texttt{1. "a road chase a dog"  =>  tipi buro dobudi} \\
\texttt{2. "a tree hold a boat"  =>  dugu nira guti} \\
\texttt{3. "a farmer chase a tree"  =>  guti buro nidi} \\
\texttt{STATUS: investigation round 3 of 4 done. Probes used: 9 of 12. 1 investigation round(s) left before the exam (10 sentences).} \\
\texttt{Send your next probes now (a PROBES list; brief THOUGHTS optional).} \\
            }
        }
    }
     \\ \\

    \theutterance \stepcounter{utterance}  
    & \multicolumn{4}{p{0.6\linewidth}}{
        \cellcolor[rgb]{0.561,0.737,0.561}{
            \makecell[{{p{\linewidth}}}]{
                \texttt{\tiny{[P1$\rangle$GM]}}
                \texttt{THOUGHTS: Now I have a clearer picture:} \\
\\ 
\texttt{Pattern: OBJECT{-}VERB{-}SUBJECT} \\
\\ 
\texttt{Verbs:} \\
\texttt{{-} chase = buro} \\
\texttt{{-} hold = nira} \\
\texttt{{-} watch = deniru} \\
\texttt{{-} hear = puga} \\
\texttt{{-} carry = geti} \\
\\ 
\texttt{Number marking:} \\
\texttt{{-} Singular nouns: dugu, tipi, guti, gutinu, nidi, bodenu, dobudi, dobudinu} \\
\texttt{{-} Plural nouns: dugunu, dobudi, bodenu (with ga), tipinu, gutinu (with ga), dobudinu (with nu)} \\
\\ 
\texttt{The verb form seems to depend on the object:} \\
\texttt{{-} buro (chase) {-} used with: dugu, tipi, guti} \\
\texttt{{-} nira (hold) {-} used with: dugu, dobudi} \\
\texttt{{-} deniru (watch) {-} used with: dugunu, nidi} \\
\texttt{{-} puga (hear) {-} used with: tipi} \\
\texttt{{-} geti (carry) {-} used with: dobudinu, pirenu} \\
\\ 
\texttt{Let me test a few more to confirm the pattern and understand the verb forms better.} \\
\\ 
\texttt{PROBES:} \\
\texttt{1. a dog hold a road} \\
\texttt{2. a boat carry a stone} \\
\texttt{3. a hill watch a dog} \\
            }
        }
    }
    & & \\ \\

    \theutterance \stepcounter{utterance}  
    & & & \multicolumn{2}{p{0.3\linewidth}}{
        \cellcolor[rgb]{0.9,0.9,0.9}{
            \makecell[{{p{\linewidth}}}]{
                \texttt{\tiny{[GM$|$GM]}}
                {'kind': 'investigation', 'round': 3, 'probes_issued': 3, 'malformed': 0, 'probes': ['a dog hold a road', 'a boat carry a stone', 'a hill watch a dog']}
            }
        }
    }
    & & \\ \\

    \theutterance \stepcounter{utterance}  
    & & & \multicolumn{4}{p{0.6\linewidth}}{
        \cellcolor[rgb]{0.9,0.9,0.9}{
            \makecell[{{p{\linewidth}}}]{
                \texttt{\tiny{[P1$\langle$GM]}}
                \texttt{TRANSLATIONS:} \\
\texttt{1. "a dog hold a road"  =>  dobudi nira tipi} \\
\texttt{2. "a boat carry a stone"  =>  bode geti dugu} \\
\texttt{3. "a hill watch a dog"  =>  tipi deniru ruditi} \\
\\ 
\texttt{FINAL EXAM. Investigation is over. Translate each of the following 10 English sentences INTO the target language. Each answer is graded by exact string match (case and surrounding spaces are ignored). Every sentence uses only words you have already seen translated.} \\
\\ 
\texttt{1. two ships hold a farmer} \\
\texttt{2. two ships chase a ship} \\
\texttt{3. a tree chase a farmer} \\
\texttt{4. three ships watch three boats} \\
\texttt{5. three hills hear two hills} \\
\texttt{6. a boat carry two hills} \\
\texttt{7. a hill hold a tree} \\
\texttt{8. two boats hold a farmer} \\
\texttt{9. a farmer hear two dogs} \\
\texttt{10. a hill watch three boats} \\
\\ 
\texttt{Respond with the answers (you may add a brief THOUGHTS line first):} \\
\texttt{ANSWERS:} \\
\texttt{1. <translation of sentence 1>} \\
\texttt{... through 10.} \\
            }
        }
    }
     \\ \\

    \theutterance \stepcounter{utterance}  
    & \multicolumn{4}{p{0.6\linewidth}}{
        \cellcolor[rgb]{0.561,0.737,0.561}{
            \makecell[{{p{\linewidth}}}]{
                \texttt{\tiny{[P1$\rangle$GM]}}
                \texttt{THOUGHTS: Let me analyze each sentence using the patterns I've discovered:} \\
\\ 
\texttt{1. "two ships hold a farmer" {-} ships (plural) = pirenu, hold = nira, farmer (singular) = nidi. Pattern: OBJECT{-}VERB{-}SUBJECT. So: pirenu nira nidi} \\
\\ 
\texttt{2. "two ships chase a ship" {-} ships (plural) = pirenu, chase = buro, ship (singular) = tipi. So: pirenu buro tipi} \\
\\ 
\texttt{3. "a tree chase a farmer" {-} tree (singular) = dugu, chase = buro, farmer (singular) = nidi. So: dugu buro nidi} \\
\\ 
\texttt{4. "three ships watch three boats" {-} ships (plural) = pirenu, watch = deniru, boats (plural) = gutinu. So: pirenu deniru gutinu} \\
\\ 
\texttt{5. "three hills hear two hills" {-} hills (plural) = ruditu, hear = puga, hills (plural) = ruditu. So: ruditu puga ruditu} \\
\\ 
\texttt{6. "a boat carry two hills" {-} boat (singular) = dugu, carry = geti, hills (plural) = ruditu. So: dugu geti ruditu} \\
\\ 
\texttt{7. "a hill hold a tree" {-} hill (singular) = ruditi, hold = nira, tree (singular) = dugu. So: ruditi nira dugu} \\
\\ 
\texttt{8. "two boats hold a farmer" {-} boats (plural) = gutinu, hold = nira, farmer (singular) = nidi. So: gutinu nira nidi} \\
\\ 
\texttt{9. "a farmer hear two dogs" {-} farmer (singular) = nidi, hear = puga, dogs (plural) = tipi. So: nidi puga tipi} \\
\\ 
\texttt{10. "a hill watch three boats" {-} hill (singular) = ruditi, watch = deniru, boats (plural) = gutinu. So: ruditi deniru gutinu} \\
\\ 
\texttt{ANSWERS:} \\
\texttt{1. pirenu nira nidi} \\
\texttt{2. pirenu buro tipi} \\
\texttt{3. dugu buro nidi} \\
\texttt{4. pirenu deniru gutinu} \\
\texttt{5. ruditu puga ruditu} \\
\texttt{6. dugu geti ruditu} \\
\texttt{7. ruditi nira dugu} \\
\texttt{8. gutinu nira nidi} \\
\texttt{9. nidi puga tipi} \\
\texttt{10. ruditi deniru gutinu} \\
            }
        }
    }
    & & \\ \\

    \theutterance \stepcounter{utterance}  
    & & & \multicolumn{2}{p{0.3\linewidth}}{
        \cellcolor[rgb]{0.9,0.9,0.9}{
            \makecell[{{p{\linewidth}}}]{
                \texttt{\tiny{[GM$|$GM]}}
                {'kind': 'exam', 'correct': 0, 'size': 10, 'results': [0, 0, 0, 0, 0, 0, 0, 0, 0, 0], 'token_similarity': 0.28809523809523807, 'answers': ['pirenu nira nidi', 'pirenu buro tipi', 'dugu buro nidi', 'pirenu deniru gutinu', 'ruditu puga ruditu', 'dugu geti ruditu', 'ruditi nira dugu', 'gutinu nira nidi', 'nidi puga tipi', 'ruditi deniru gutinu'], 'targets': ['nidi nira pirenu nu', 'pire buro pirenu nu', 'nidi buro guti', 'dugunu ga deniru pirenu ga', 'ruditinu nu puga ruditinu ga', 'ruditinu nu geti dugu', 'guti nira ruditi', 'nidi nira dugunu nu', 'tipinu nu puga nidi', 'dugunu ga deniru ruditi']}
            }
        }
    }
    & & \\ \\

\end{supertabular}
}

\end{document}
